\documentclass[11pt,onecolumn]{IEEEtran}
\usepackage{amsmath,amssymb,amsthm}
\usepackage[hidelinks]{hyperref}

\newtheorem{theorem}{Theorem}
\newtheorem{proposition}[theorem]{Proposition}
\newtheorem{lemma}[theorem]{Lemma}
\newtheorem{corollary}[theorem]{Corollary}
\newtheorem{remark}[theorem]{Remark}

\newcommand{\E}{\mathbb{E}}
\newcommand{\R}{\mathbb{R}}
\newcommand{\tr}{\operatorname{tr}}
\newcommand{\PC}{P_{\mathcal C}}
\newcommand{\cR}{\mathcal R}
\newcommand{\nup}{\nu_{+}}
\newcommand{\num}{\nu_{-}}
\newcommand{\epp}{\varepsilon_{+}}
\newcommand{\emm}{\varepsilon_{-}}
\newcommand{\lmin}{\lambda_{\min}}
\newcommand{\lmax}{\lambda_{\max}}

\begin{document}

\title{Learn Then Serve: Recovering a Consumer's Read Operator and Paying for It\\ (T1, working note)}
\author{Andrew~H.~Bond%
\thanks{Working note of the Observation Theory program, v4.1, 2026-11-18. Drafted by
Claude (Opus 5.5) at the author's request. The author read versions 2 and 3 and asked
for the changes of version 4: the count of consumer queries qualified by the error of
finite differences (Proposition~\ref{prop:fd}, Corollary~\ref{cor:n}); the Lipschitz
statements made for the weighted map $G^{1/2}\mathcal C$ (Section~\ref{sec:setting},
Remark~\ref{rem:restrictive}); the calibration protocol stated and its message counted
(Assumption (A3), Proposition~\ref{prop:calib}); and the test-channel results made the
centre of the note. The author also proposed the route of Theorem~\ref{thm:codes},
through the finite-alphabet coding theorem of \cite[App.~A]{part1}. Version 2 applied the
fresh-context audit VI-22 and version 3 the audit VI-23. The audit VI-24 of version 4
found no false mathematical step, and version 4.1 applies its findings: one gap in
what Proposition~\ref{prop:calib} calls leakage, one overstatement in
Remark~\ref{rem:restrictive}, and wording. Numerical checks: \texttt{verify\_t1\_v4.py} in this folder. Ledger class
\texttt{[exploratory]}. Numbering, version 3 to version 4: Lemma~1 is Lemma~1;
Corollary~5 is Corollary~3; Theorem~6 is Theorem~5; Remark~7 is Remark~6; Corollary~8
is Corollary~8; Proposition~2 is Proposition~10; Theorem~3 is Theorem~11;
Proposition~4 is Proposition~12; Remarks~9--12 are Remarks~13--16; Remark~13 is split
into Proposition~2 (finite differences) and Remark~17 (noise); Remarks~14--16 are
Remarks~18--20.}}
\maketitle

\begin{abstract}
A consumer reads a source through a map $\mathcal C$ and an output metric $G$; its
loss is the squared distance between the images, $\|G^{1/2}(\mathcal C(\hat x)-\mathcal C(x))\|^2$,
and its local read metric $\PC(x)=J(x)^{\mathsf T}GJ(x)$ varies with the point $x$.
The coding theorems of the rate--leakage--distortion series
\cite{part1,part2,part3,part4} price a description once a fixed per-symbol quadratic
weight is supplied. This note connects those prices to a weight learned from the
consumer. The read metrics at $m$ probe points have a mean $\hat Q$ with
$(1-\emm)Q\preceq\hat Q\preceq(1+\epp)Q$ with probability at least $1-\delta$, where $Q$
is their expectation, $\varepsilon_\pm=\sqrt{2\nup\num L/m}+2\nu_\pm L/(3m)$,
$L=\ln(2p/\delta)$, and $\nu_\pm$ bound the read metric relative to $Q$; the lower half
needs no prior constant. Central differences of step $h$ add a deterministic error of
order $h^2$ (Proposition~\ref{prop:fd}), and a decoder that does not probe the consumer
can be sent the weight once, in $O(p^2\log(p\kappa/\varepsilon_q))$ bits that carry no
information about the source (Proposition~\ref{prop:calib}). For a Gaussian vector source and any consumer
with three bounded derivatives, the Gaussian test channel's consumer loss is
$\tr(Q\Delta)$ up to an explicit remainder of order $(\tr\Delta)^2$, with no
sandwich assumed (Theorem~\ref{thm:second}), and, by the coding theorem of
\cite{part1}, codes achieve the test channel's rate, leakage, and consumer loss
together (Theorem~\ref{thm:codes}). So a learner that knows derivative bounds, the
source law, and an upper constant for the read metric can certify that the series'
frontier for $\hat Q$, at a budget $\hat D$ slightly below $D$, is achievable at consumer
loss $D$ by codes built for the consumer, within
$\bar\tau=\tau_\varepsilon+O(D)$ bits per symbol of the frontier for the true weight,
$\tau_\varepsilon=\tfrac p2\log\frac{1+\epp}{1-\emm}$ (Corollary~\ref{cor:cert}). For
consumers under a relative quadratic sandwich known in advance, whose weighted map is
then bi-Lipschitz, a converse holds as well: in every setting of the series the codes
for $\hat Q$ lie within $\tau_\nu+\tau_\varepsilon$ of the outer bound
(Theorem~\ref{thm:T1}), and the learning part is sharp for an encoder that knows the
weight only through the confidence statement (Proposition~\ref{prop:sharp}). Open:
codes in the process settings, one code that serves every consumer of a class, and a
converse without the sandwich.
\end{abstract}

\section{Setting}\label{sec:setting}

\textbf{The coding side.} Fix any setting of \cite{part1,part2,part3,part4}
in which a source of vector symbols $T_t\in\R^p$ is described to a decoder
under the per-symbol quadratic distortion
$\tfrac1N\sum_{t\le N}\E(T_t-\hat T_t)^{\mathsf T}W(T_t-\hat T_t)$ with a weight
$W\succ0$ (blocklength $N$), while a holder with a fixed observation model
measures the leakage. Write $\cR_W(D)$ for the operational region of pairs
of rate and leakage at the distortion $D$: the closure of the achievable
pairs, which the theorems of the series identify with the set of pairs
above every supporting line $\alpha R+(1-\alpha)L\ge V_W(\alpha,D)$,
$\alpha\in[0,1]$, for the frontier cost $V_W$ of the setting (Gaussian
vectors \cite{part1}; stationary processes \cite{part2}; locally stationary,
cyclostationary, and scalar unstable sources \cite{part3}; unstable vector,
rational-filter, unit-circle scalar, and ergodically driven sources
\cite{part4}). At and beyond the setting's maximal distortion the region is
the nonnegative quadrant, or its translate by the setting's shift, and we
read every region statement below with that convention, so that no range
condition on $D$ is needed. We use three properties of $V_W$, each a
one-line consequence of its definition as a minimum over error covariances
$0\prec\Delta\preceq\mathsf S$ of $\alpha r(\Delta)+(1-\alpha)\ell(\Delta)$ under
the budget $\tr((I\otimes W)\Delta)\le ND$, passed to the limits that define
the process settings:
\begin{itemize}
\item[(P1)] \emph{Scaling:} $V_{cW}(\alpha,D)=V_W(\alpha,D/c)$ for $c>0$.
\item[(P2)] \emph{Monotonicity:} $W\preceq W'$ implies $V_W\le V_{W'}$.
\item[(P3)] \emph{Budget inflation:} $V_W(\alpha,D/c)\le V_W(\alpha,D)+\tfrac p2\log c$
for $c\ge1$.
\end{itemize}
(P1) holds because the two budgets define the same feasible set. (P2)
holds because a larger weight shrinks the feasible set. (P3): if $\Delta$ is
feasible at $D$ then $\Delta/c$ is feasible at $D/c$; its rate is larger by
exactly $\tfrac{pN}2\log c$ per block of $N$ symbols, and its leakage
$\tfrac12\log\det(c\Delta^{-1}+\mathsf J)-\tfrac12\log\det(\mathsf S^{-1}+\mathsf J)$
is larger by at most the same, because
$c\Delta^{-1}+\mathsf J\preceq c(\Delta^{-1}+\mathsf J)$. This is the inequality of
\cite[Lemma~11(b),(c)]{part2}, which \cite[Sec.~II(iii)]{part4} re-derives
for any information matrix. In the settings of \cite{part3,part4} with
shifts, such as $(\rho_\star,\rho_{\mathrm{hid}})$ for unstable poles, the shift
does not depend on $W$ and the three properties hold for the shifted
frontier. Because a closed, convex, upward-closed set is the intersection of
its supporting half-planes \cite[proof of Theorem~16]{part2}, (P1)--(P3)
give, for the regions,
\[
\cR_{cW}(D)=\cR_W(D/c),\qquad W\preceq W'\Rightarrow\cR_{W'}(D)\subseteq\cR_W(D),\qquad
\cR_W(D/c)\supseteq\cR_W(D)+\tfrac p2\log c\,(1,1)\ (c\ge1).
\]
Logarithms are to base 2 unless written $\ln$. $V_W(\alpha,\cdot)$ is nonincreasing,
and for $D'>D$ (P3) gives $V_W(\alpha,D)\le V_W(\alpha,D')+\tfrac p2\log(D'/D)$; so
$V_W(\alpha,\cdot)$ is right-continuous, and $\bigcap_{D'>D}\cR_W(D')=\cR_W(D)$.

\textbf{The consumer side.} A consumer $\mathcal C:\R^p\to\R^k$ acts on each
symbol, with a constant output metric $G\succeq0$ and the loss
\[
\ell_{\mathcal C}(x,\hat x)=\|\mathcal C(\hat x)-\mathcal C(x)\|_G^2=\|g(\hat x)-g(x)\|^2,
\qquad g=G^{1/2}\mathcal C .
\]
The loss sees $\mathcal C$ only through the weighted map $g$, and every regularity
statement below (Lipschitz constants, derivative bounds, injectivity) is a statement
about $g$. It transfers to $\mathcal C$ when $G\succ0$, and not otherwise: the part of
$\mathcal C$ in the kernel of $G$ is unrestricted (Remark~\ref{rem:restrictive}). Where $g$ is differentiable, the local read metric is
$\PC(x)=J_g(x)^{\mathsf T}J_g(x)$, $J_g=\partial g/\partial x$, which is
$J^{\mathsf T}GJ$, $J=\partial\mathcal C/\partial x$, where $\mathcal C$ is differentiable
\cite{otieee}. When $g$ is $C^j$ write $M_j=\sup_x\|D^jg(x)\|$ for $j=1,2,3$
(operator norms of the $j$-linear maps, Euclidean on $\R^p$). Let $\pi$ be a probe
distribution on $\R^p$ and $Q=\E_\pi\PC(X)$.

\textbf{Assumption (A2), probing.} The learner, which sits at the encoder, draws
$x_1,\dots,x_m$ i.i.d.\ from $\pi$, obtains the read metric at each point, and forms
$\hat Q=\tfrac1m\sum_i\PC(x_i)$. Exact read metrics need the Jacobian of $g$; from
consumer queries, central differences give it at $2p$ queries per point, with the
error that Proposition~\ref{prop:fd} bounds.

\textbf{Assumption (A3), calibration.} Coding starts after a calibration phase. Encoder
and decoder use the code that the setting's achievability theorem gives for the weight
$\hat Q$, which both construct, by a deterministic selection rule agreed in advance among
the codes the theorem asserts to exist, from $\hat Q$, the constants of the confidence
statement (and, for Corollary~\ref{cor:cert}, the derivative bounds $M_1,M_2,M_3$, which
enter its budget), and the source and holder models that the theorem assumes known. The
decoder obtains $\hat Q$ in one of two ways.
\begin{itemize}
\item[(A3-i)] \emph{Shared calibration.} It holds the same probe points (for instance
through a common seed) and query access to the same consumer, and computes the same
$\hat Q$. Nothing is communicated; the decoder's own queries are counted with the
learner's.
\item[(A3-ii)] \emph{A calibration message.} The learner sends a quantized weight
$\tilde Q$ once, and encoder and decoder both use $\tilde Q$. Proposition~\ref{prop:calib}
counts its bits and what it leaks.
\end{itemize}
Three resources are then bounded separately: consumer queries ($2pm$ at the learner,
once, and as many again at the decoder under (A3-i)), calibration bits (none under
(A3-i), at most $B$ once under (A3-ii)), and the rate and leakage of the code per symbol.
Over a session of $K$ symbols the one-time costs are $2pm/K$ queries ($4pm/K$ under
(A3-i)) and $B/K$ bits per symbol. This note does not optimize a combined budget; a price of a
query in bits is not defined here.

\section{Learning the weight}

\begin{lemma}[Estimating the weight]\label{lem:est}
Suppose (A2), $Q\succ0$, and $(1-\num)Q\preceq\PC(x)\preceq(1+\nup)Q$ on the support
of $\pi$, with $\nup\ge0$ and $\num\in[0,1]$. Then with probability at least
$1-\delta$,
\[
(1-\emm)Q\preceq\hat Q\preceq(1+\epp)Q,\qquad
\varepsilon_\pm=\sqrt{\frac{2\nup\num L}m}+\frac{2\nu_\pm L}{3m},\quad L=\ln\frac{2p}\delta .
\]
\end{lemma}

\begin{proof}
Let $X_i=Q^{-1/2}\PC(x_i)Q^{-1/2}-I$. The $X_i$ are independent, symmetric, and
centred, with $-\num I\preceq X_i\preceq\nup I$. The polynomial
$f(t)=(\nup-t)(t+\num)$ is nonnegative on $[-\num,\nup]$, which contains the
spectrum of $X_i$, so $f(X_i)\succeq0$, that is,
$X_i^2\preceq(\nup-\num)X_i+\nup\num I$; taking expectations,
$\E X_i^2\preceq\nup\num I$ (the matrix form of the Bhatia--Davis inequality
\cite{bhatiadavis}). So $\|\sum_i\E X_i^2\|\le m\nup\num$. The matrix
Bernstein inequality \cite[Theorem~1.4]{tropp2012} needs only an upper bound
$\lmax(X_i)\le R$; applied to $\{X_i\}$ with $R=\nup$ and to $\{-X_i\}$ with
$R=\num$, it gives
\[
\Pr\Bigl\{\lmax\Bigl(\sum_iX_i\Bigr)\ge mt\Bigr\}\le p\,e^{-\frac{mt^2/2}{\nup\num+\nup t/3}},\qquad
\Pr\Bigl\{\lmax\Bigl(-\sum_iX_i\Bigr)\ge mt\Bigr\}\le p\,e^{-\frac{mt^2/2}{\nup\num+\num t/3}}.
\]
The exponents equal $-L$ at
$t_\pm=\frac{\nu_\pm L}{3m}+\sqrt{\bigl(\frac{\nu_\pm L}{3m}\bigr)^2+\frac{2\nup\num L}m}\le\varepsilon_\pm$,
respectively, so each tail has probability at most $pe^{-L}=\delta/2$ at $t=\varepsilon_\pm$. Outside both,
$-\emm I\preceq\tfrac1m\sum_iX_i=Q^{-1/2}\hat QQ^{-1/2}-I\preceq\epp I$.
\end{proof}

Write $\nu=\max\{\num,\nup\}$. Since $\nup\num\le\nu^2$, both $\varepsilon_\pm$ are at
most the constant $\nu(\sqrt{2L/m}+2L/(3m))$ of version~2. The variance bound is
attained by the two-point law $X=\nup$ with probability $\num/(\nup+\num)$ and $X=-\num$
otherwise. When either constant is zero the square-root term vanishes and the
bound is of order $1/m$ (in fact a read metric that never falls below its mean is
constant, and the error is zero). Because $\PC(x)\succeq0$, the lower half of the
hypothesis holds with $\num=1$ for every consumer, so only an upper constant $\nup$ must
be known in advance. On the registration draft's consumer $x+0.2\sin x$, $m=32$,
$\delta=0.1$, the constants are $\epp=0.186$ and $\emm=0.208$, against $0.291$ in
version~2; the exact roots $t_\pm$ of the proof, $0.177$ and $0.189$, may be used in
their place. We call
\[
\tau_\varepsilon=\frac p2\log\frac{1+\epp}{1-\emm}
\]
the learning price: it is what knowing $Q$ only through the event of
Lemma~\ref{lem:est} costs, in bits per symbol, in Corollary~\ref{cor:cert} and
Theorem~\ref{thm:T1}, and it is sharp in the sense of Proposition~\ref{prop:sharp}.

\begin{proposition}[Finite differences]\label{prop:fd}
Let $g$ be $C^2$, $x\in\R^p$, and $h>0$. Let $\hat J$ have the columns
$(\mathcal C(x+he_j)-\mathcal C(x-he_j))/(2h)$, $j=1,\dots,p$, from $2p$ noiseless consumer
queries, and $\hat P(x)=\hat J^{\mathsf T}G\hat J$. Then
\[
\|\hat P(x)-\PC(x)\|\le 2M_1\sqrt p\,\beta+p\beta^2,\qquad
\beta=\min\Bigl\{\frac{M_2h}2,\ \frac{M_3h^2}6\Bigr\},
\]
the second entry of the minimum when $g$ is $C^3$. Consequently, if
$0<\lambda\le\lmin(Q)$, the mean $\hat Q_h$ of $\hat P$ over the probe points of (A2)
satisfies, on the event of Lemma~\ref{lem:est},
\[
(1-\emm-\varepsilon_h)Q\preceq\hat Q_h\preceq(1+\epp+\varepsilon_h)Q,\qquad
\varepsilon_h=\frac{2M_1\sqrt p\,\beta+p\beta^2}{\lambda}.
\]
\end{proposition}

\begin{proof}
The $j$th column of $G^{1/2}\hat J$ is $(g(x+he_j)-g(x-he_j))/(2h)$. By Taylor's
theorem with integral remainder,
$g(x\pm he_j)=g(x)\pm hJ_g(x)e_j+\int_0^h(h-t)D^2g(x\pm te_j)[e_j,e_j]\,dt$, and the
two remainders have norm at most $M_2h^2/2$ each, so the column differs from
$J_g(x)e_j$ by at most $M_2h/2$. When $g$ is $C^3$, expanding one order further, the
second-order terms cancel and the two remainders
$\pm\int_0^h\tfrac12(h-t)^2D^3g(x\pm te_j)[e_j,e_j,e_j]\,dt$ have norm at most $M_3h^3/6$
each, so the column error is at most $M_3h^2/6$. So $B=G^{1/2}\hat J-J_g(x)$ has
columns of norm at most $\beta$, and $\|B\|\le\|B\|_F\le\sqrt p\,\beta$. Since
$\hat P-\PC=J_g^{\mathsf T}B+B^{\mathsf T}J_g+B^{\mathsf T}B$ and $\|J_g\|\le M_1$, the
first bound follows. For the second, $|u^{\mathsf T}(\hat P(x_i)-\PC(x_i))u|\le
\varepsilon_h\lambda\|u\|^2\le\varepsilon_hu^{\mathsf T}Qu$ for every $u$ and every $i$, so
$\hat Q_h-\hat Q$ lies between $-\varepsilon_hQ$ and $\varepsilon_hQ$; add the event of
Lemma~\ref{lem:est}.
\end{proof}

The perturbation is deterministic, so it adds to $\varepsilon_\pm$ rather than to the
variance, and it vanishes as $h\to0$ for a noiseless consumer. Output noise changes
this (Remark~\ref{rem:fd}).

\begin{corollary}[Probes for a target price]\label{cor:n}
Let $0<\eta\le p$.
\begin{itemize}
\item[(i)] \emph{Exact read metrics.} The learning price is $\tau_\varepsilon\le\eta$
bits per symbol if
\[
m\ \ge\ L\max\Bigl\{\frac{200}9\,\frac{\nup\num p^2}{\eta^2},\ \frac{20}9\,\frac{\nu p}\eta\Bigr\}.
\]
The probing then costs $m$ read metrics. Their cost in consumer queries depends on how
the Jacobians are obtained, and is not counted here.
\item[(ii)] \emph{Central differences.} With $g$ of class $C^3$ and $0<\lambda\le\lmin(Q)$,
the learning price computed with $\varepsilon_\pm+\varepsilon_h$ in place of
$\varepsilon_\pm$ is at most $\eta$ if
\[
m\ \ge\ L\max\Bigl\{50\,\frac{\nup\num p^2}{\eta^2},\ \frac{10}3\,\frac{\nu p}\eta\Bigr\}
\quad\text{and}\quad
h^2\ \le\ \frac6{M_3}\min\Bigl\{\frac{\lambda\eta}{20M_1p^{3/2}},\ \Bigl(\frac{\lambda\eta}{10p^2}\Bigr)^{1/2}\Bigr\}.
\]
The probing then costs $2pm$ noiseless consumer queries, once; served over $K$ symbols
it adds $2pm/K$ queries per symbol.
\end{itemize}
\end{corollary}

\begin{proof}
Let $t=\tanh(\eta\ln2/p)$. If the two relative errors are at most $t$, the price is at
most
\[
\tfrac p2\log\tfrac{1+t}{1-t}=p\,\mathrm{artanh}(t)/\ln2=\eta .
\]
Both
$\varepsilon_\pm\le\sqrt{2\nup\num L/m}+2\nu L/(3m)$. Since $\tanh$ is concave on
$[0,\infty)$ and $\tanh(\ln2)=3/5$, $t\ge\tfrac35\eta/p$ for $\eta\le p$.
(i) Each of the two terms is at most $t/2$ when $m\ge8\nup\num L/t^2$ and
$m\ge4\nu L/(3t)$, and substituting the bound on $t$ gives the two terms.
(ii) Each of the two terms, and $\varepsilon_h$, is at most $t/3$. For the terms
this holds when $m\ge18\nup\num L/t^2$ and $m\ge2\nu L/t$, which the stated $m$
implies. For $\varepsilon_h$, it is at most $\eta/(5p)\le t/3$ when both
$2M_1\sqrt p\,\beta$ and $p\beta^2$ are at most $\lambda\eta/(10p)$, which the stated
$h$ implies through $\beta\le M_3h^2/6$.
\end{proof}

Without an a priori lower constant, $\num=1$ and the first terms read
$\tfrac{200}9\nup p^2L/\eta^2$ and $50\nup p^2L/\eta^2$.

\begin{proposition}[The calibration message]\label{prop:calib}
Let $0<\lambda\le\lmin(Q)$, $\Lambda\ge\lmax(Q)$, $0<\varepsilon_q<1$, and the constants
$e_\pm$ below be known to both ends, and $\kappa=\Lambda/\lambda$. Let $\hat Q$ be the learner's weight,
with $(1-e_-)Q\preceq\hat Q\preceq(1+e_+)Q$ on an event, $e_-<1$ (for instance
$e_\pm=\varepsilon_\pm$ by Lemma~\ref{lem:est}, or $\varepsilon_\pm+\varepsilon_h$ by
Proposition~\ref{prop:fd}). Round each entry of $\hat Q$ on or above the diagonal to the
nearest point of the grid $\{js:|j|\le N_q\}$, with
\[
s=\frac{2(1-e_-)\lambda\varepsilon_q}p,\qquad
N_q=\Bigl\lceil\frac{p\kappa(1+e_+)}{2(1-e_-)\varepsilon_q}\Bigr\rceil,
\]
and let $\tilde Q$ be the symmetric matrix of the rounded entries. Then $\tilde Q$ is
described by
\[
B=\frac{p(p+1)}2\bigl\lceil\log(2N_q+1)\bigr\rceil
\ =\ O\Bigl(p^2\log\frac{p\kappa(1+e_+)}{(1-e_-)\varepsilon_q}\Bigr)
\]
bits, and on the event
\[
(1-e_-)(1-\varepsilon_q)Q\preceq\tilde Q\preceq(1+e_+)(1+\varepsilon_q)Q .
\]
Sent once before $K$ symbols are coded, the message adds $B/K$ bits per symbol to the
rate. It carries no information about the source: for the code's message $M$ and the
holder's observations $S^K$, $I(T^K;\tilde Q,M\mid S^K)=I(T^K;M\mid S^K,\tilde Q)$. So
when the realized calibration is part of the code and known, as a code is, to every
party, the leakage is that of the code designed for $\tilde Q$. In the series' form
$\tfrac1KH(\tilde Q,M\mid S^K)$, which also counts the message's own entropy, it adds at
most $B/K$ bits per symbol.
\end{proposition}

\begin{proof}
On the event every entry of $\hat Q$ has absolute value at most
$\lmax(\hat Q)\le(1+e_+)\Lambda\le N_qs$, so it lies in the range of the grid and is
moved by at most $s/2$. The error $E=\tilde Q-\hat Q$ is symmetric with entries of
absolute value at most $s/2$, so $\|E\|\le\|E\|_F\le ps/2=(1-e_-)\lambda\varepsilon_q
\le\varepsilon_q\lmin(\hat Q)$, and $-\varepsilon_q\hat Q\preceq E\preceq\varepsilon_q\hat Q$;
compose with the event. The grid has $2N_q+1$ points and there are $p(p+1)/2$ entries.
For the leakage, $\tilde Q$ is a function of the probe points and the consumer, which
are independent of the source and of the holder's observations, so
$I(T^K;\tilde Q\mid S^K)=0$ and $I(T^K;\tilde Q,M\mid S^K)=I(T^K;M\mid S^K,\tilde Q)$. Given
$\tilde Q=q$ the code is deterministic, so $H(M\mid S^K,\tilde Q=q)=I(T^K;M\mid S^K,\tilde Q=q)$,
the leakage of the code designed for $q$. Unconditionally,
$H(\tilde Q,M\mid S^K)=H(\tilde Q)+H(M\mid S^K,\tilde Q)$ because $\tilde Q$ is independent of
$S^K$, and $H(\tilde Q)\le B$. The same holds under (A3-i) with the probe points in place
of $\tilde Q$: if the holder did not know them, the entropy form could exceed the
designed leakage, while the conditional information could not.
\end{proof}

All results below are stated for the weight $\hat Q$ of Lemma~\ref{lem:est}. They hold
for the weight that encoder and decoder actually share: for $\hat Q_h$ with
$\varepsilon_\pm+\varepsilon_h$ in place of $\varepsilon_\pm$, and for $\tilde Q$ with
$1+\epp$ replaced by $(1+e_+)(1+\varepsilon_q)$ and $1-\emm$ by $(1-e_-)(1-\varepsilon_q)$,
which adds $\tfrac p2\log\frac{1+\varepsilon_q}{1-\varepsilon_q}$ to the learning price.

\section{Any smooth consumer: what the test channel pays and what codes achieve}

For a consumer that is not sandwiched by a quadratic form there is no converse in terms
of a weight (Remark~\ref{rem:second}), but the frontiers of the series are attained by
Gaussian test channels, whose errors are nearly independent of the source. For them
the consumer loss is $\tr(Q\Delta)$ to second order in the distortion, for any consumer
with three bounded derivatives (Theorem~\ref{thm:second}); in the memoryless setting
codes achieve what the test channel achieves (Theorem~\ref{thm:codes}); and the two
together certify a learned design (Corollary~\ref{cor:cert}).

Take the setting of \cite{part1} with the decoder reproducing the source:
$T\sim N(0,\Sigma)$ with $\Sigma\succ0$, a holder $S=HT+U$. For
$0\prec\Delta\preceq\Sigma$, the Gaussian test channel with error covariance $\Delta$
has $T-\hat T$ independent of $\hat T$ with covariance $\Delta$; equivalently
\[
\hat T=(I-\Delta\Sigma^{-1})T+V,\qquad V\sim N(0,\Delta-\Delta\Sigma^{-1}\Delta)\ \text{independent of }(T,U),
\]
and the frontier points of \cite{part1} are the rate and leakage
$(r(\Delta),\ell(\Delta))=(I(T;\hat T),I(T;\hat T\mid S))$ of these channels at the
minimizing $\Delta$. For a consumer whose weighted map $g$ is $C^3$ with $M_1,M_2,M_3$
finite, take the probe law to be the source law, $\pi=N(0,\Sigma)$, so that
$Q=\E\PC(T)$, and let $\omega=\sup_x\|\PC(x)-Q\|$. Since
$0\preceq\PC\preceq M_1^2I$, $\omega\le\max\{M_1^2-\lmin(Q),\|Q\|\}\le M_1^2$. With
$s=\tr\Delta$ and $\theta=\tr(\Delta\Sigma^{-1}\Delta)\le s^2/\lmin(\Sigma)$, let
\[
\rho_b(\Delta)=\bigl(\tfrac74M_2^2+2M_1M_3\bigr)s^2+\tfrac{\sqrt5}2M_2M_3s^{5/2}+\tfrac5{12}M_3^2s^3,
\]
\[
\rho_f(\Delta)=\omega\theta+M_1M_2\bigl(3\theta^{1/2}s+3^{3/4}\theta^{3/2}\bigr)
+\bigl(\tfrac34M_2^2+M_1M_3\bigr)s^2+\tfrac{\sqrt5}2M_2M_3s^{5/2}+\tfrac5{12}M_3^2s^3 .
\]

\begin{theorem}[The test channel is exact to second order]\label{thm:second}
For every $0\prec\Delta\preceq\Sigma$,
\[
\bigl|\E\,\ell_{\mathcal C}(T,\hat T)-\tr(Q\Delta)\bigr|\ \le\ \rho(\Delta)=\min\{\rho_b(\Delta),\rho_f(\Delta)\}.
\]
Both bounds are $O(s^2)$, and neither dominates the other. For a linear consumer
$\rho=0$: the identity $\E\,\ell_{\mathcal C}=\tr(Q\Delta)$ is then exact.
\end{theorem}

\begin{proof}
\emph{Taylor.} For fixed $t$,
$g(t+\delta)-g(t)=J_g(t)\delta+\tfrac12D^2g(t)[\delta,\delta]+r$
with $\|r\|\le M_3\|\delta\|^3/6$. Expanding the squared norm,
\[
\ell_{\mathcal C}(t,t+\delta)=\delta^{\mathsf T}\PC(t)\delta+\kappa_t(\delta)+R,\qquad
\kappa_t(\delta)=\bigl\langle J_g(t)\delta,\,D^2g(t)[\delta,\delta]\bigr\rangle,
\]
\[
|R|\le\bigl(\tfrac14M_2^2+\tfrac13M_1M_3\bigr)\|\delta\|^4+\tfrac16M_2M_3\|\delta\|^5+\tfrac1{36}M_3^2\|\delta\|^6 .
\]

For $\delta\sim N(0,\Delta)$, Jensen's inequality on the eigen-decomposition gives
$\E\|\delta\|^{2k}\le(2k-1)!!\,s^k$, so $\E\|\delta\|^4\le3s^2$,
$\E\|\delta\|^6\le15s^3$, $\E\|\delta\|^5\le(3\cdot15)^{1/2}s^{5/2}$, and $\E|R|$ is at
most $(\tfrac34M_2^2+M_1M_3)s^2+\tfrac{\sqrt5}2M_2M_3s^{5/2}+\tfrac5{12}M_3^2s^3$.

\emph{The bound $\rho_b$: expand about $\hat T$.} The loss is symmetric in its two
arguments, and $Z=T-\hat T\sim N(0,\Delta)$ is independent of $\hat T$. Expand with
$t=\hat T$ and $\delta=Z$. The cubic term is a homogeneous cubic form in $Z$ whose
coefficients depend on $\hat T$ only, so its mean is zero. The quadratic term has mean
$\tr(\E\PC(\hat T)\,\Delta)$. Since $T=\hat T+Z$, $Q=\E\PC(\hat T+Z)$; by Taylor's theorem
with integral remainder, the first-order term $D\PC(\hat T)[Z]$ has mean zero, and
$D^2\PC[z,z]=2(D^2g[z,\cdot])^{\mathsf T}(D^2g[z,\cdot])
+(D^3g[z,z,\cdot])^{\mathsf T}J_g+J_g^{\mathsf T}(D^3g[z,z,\cdot])$ has norm at
most $2(M_2^2+M_1M_3)\|z\|^2$, so $\|Q-\E\PC(\hat T)\|\le(M_2^2+M_1M_3)s$ and
$|\tr((Q-\E\PC(\hat T))\Delta)|\le(M_2^2+M_1M_3)s^2$. Adding $\E|R|$ gives $\rho_b$.

\emph{The bound $\rho_f$: expand about $T$.} Write $\delta=\hat T-T=u+V$ with
$u=-\Delta\Sigma^{-1}T$, so that $u$ and $V$ are independent centred Gaussian vectors
with $\E\|u\|^2=\theta$ and $\E\|V\|^2=s-\theta\le s$. The cross term
$2\E u^{\mathsf T}\PC(T)V$ vanishes, and
$\E V^{\mathsf T}\PC(T)V=\tr(Q(\Delta-\Delta\Sigma^{-1}\Delta))$ by independence, while
$\E u^{\mathsf T}Qu=\tr(Q\Delta\Sigma^{-1}\Delta)$. So
$\E\delta^{\mathsf T}\PC(T)\delta=\tr(Q\Delta)+\E u^{\mathsf T}(\PC(T)-Q)u$, and the last
term is at most $\omega\theta$ in absolute value. Given $T$, the vector $V$ is centred, Gaussian, and symmetric,
so the parts of the trilinear form $\kappa_T(u+V)$ of degree one and three in $V$ have
conditional mean zero. The remaining parts are $\kappa_T(u)$ and three terms with one
$u$ and two $V$'s, each at most $M_1M_2\|u\|\|V\|^2$. With
$\E\|u\|^3\le(\E\|u\|^4)^{3/4}\le(3\theta^2)^{3/4}$ and independence,
$|\E\kappa_T(\delta)|\le M_1M_2(3^{3/4}\theta^{3/2}+3\theta^{1/2}s)$. Adding $\E|R|$ gives
$\rho_f$.
\end{proof}

\begin{remark}[What Theorem~\ref{thm:second} does and does not say]\label{rem:second}
It assumes no sandwich. The consumer may be non-injective, saturating, or have a
singular read metric at some points: the circle consumer $x\mapsto(\cos x,\sin x)$,
which fails the sandwich (A1a) of Section~\ref{sec:sandwich}, has $\E\ell_{\mathcal C}=2(1-e^{-\Delta/2})=\Delta-\Delta^2/4+\dots$.
The relative excess is $O(s)$, so the factor $1+\nup$ of Theorem~\ref{thm:T1} is not
paid: for the registration draft's consumer $x+0.2\sin x$ (where $1+\nup=1.233$) and its
design at $D=0.3\tr(Q\Sigma)$, the measured ratio $\E\ell_{\mathcal C}/\tr(Q\Delta)$ is $0.998$
(\texttt{verify\_t1\_v4.py}, check C5). The audit of version~1 (ledger VI-22)
reported a Monte Carlo loss of $0.301$ for $x+0.05\sin x$ at $D=0.3$; for the instance
of these checks the exact loss of that design is $0.29999$ (audit VI-23), so the
excess was sampling noise, if the two instances agree (VI-22 does not record its own).
The constants are not tight: over the checks $\rho$ exceeds the measured excess by a
factor of at least $30$, and the audit found cases within a factor of $12$ for
$\rho_f$ (the circle consumer with $\Sigma=10$). For consumers with four derivatives a
formal expansion, derived and checked numerically in audit VI-23 but not proved here,
gives the excess as $c_2(\Delta)+o(s^2)$ with, for $G=I$ and $H_i$ the Hessian of the
$i$th output, $c_2=\E\sum_i[\tfrac14(\tr H_i\Delta)^2-\tfrac12\tr(H_i\Delta H_i\Delta)]$;
in one dimension $c_2=-\tfrac14\E[\mathcal C''(T)^2]\Delta^2<0$, so the test channel
loses slightly less than $\tr(Q\Delta)$, but for $p\ge3$ the sign can be positive. Three
limits. The leading term needs the probe law to be the source law; for another $\pi$ it
is $\tr(\E_{\mathrm{source}}\PC\,\Delta)$. The statement is about the single-letter test
channel; Theorem~\ref{thm:codes} carries it to codes in this memoryless setting, and in
the process settings of \cite{part2,part3,part4} the transfer is open. And there is no
converse: for a non-injective consumer a code can exploit the non-injectivity (the
circle consumer needs only $x$ modulo $2\pi$), so $\cR_Q(D)$ is then not a lower bound.
The best code for a locally quadratic loss itself, rather than for a constant weight,
is the subject of the high-resolution theory of non-difference distortion measures
\cite{linderzamir}, which this note has not read at first hand and does not use.
\end{remark}

\begin{theorem}[From test channels to codes]\label{thm:codes}
In the setting of \cite{part1} with the decoder reproducing $T$, let $g$ be Lipschitz
with constant $M_1$. Let $\hat T$ be the output of a channel from $T$, with $\hat T-T-S$,
$\E\|\hat T\|^2<\infty$, and $I(T;\hat T)<\infty$. For every $\epsilon>0$ and every large
$n$ there is a deterministic code of blocklength $n$ with rate at most
$I(T;\hat T)+\epsilon$, leakage $\tfrac1nH(M\mid S^n)$ at most $I(T;\hat T\mid S)+\epsilon$,
and consumer loss per symbol $\tfrac1n\sum_i\E\,\ell_{\mathcal C}(T_i,\hat T_i)$ at most
$\E\,\ell_{\mathcal C}(T,\hat T)+\epsilon$. In particular the pair $(r(\Delta),\ell(\Delta))$
of the Gaussian test channel is achievable at that channel's consumer loss. The code
depends on the consumer.
\end{theorem}

\begin{proof}
Follow the achievability part of the proof of \cite[Theorem~3]{part1}, in its
Appendix~A, with the described variable equal to $T$ and its surrogate distortion
replaced by
\[
d_k(t,\hat x)=\E\bigl[\ell_{\mathcal C}(T,\hat x)\bigm|T_k=t\bigr]
\]
for the quantized source $T_k=q_k(T)$ and a reproduction letter $\hat x$. The
quantization of the side information (Step~A), the bounds on rate and leakage in
Step~B, and the diagonal argument in the first part of Step~C do not involve the
distortion. The rest of Step~C, which identifies the region, is not needed here, and the
proof's opening line, which derives $\E\|\hat Y\|^2<\infty$ from a finite quadratic
distortion and $W\succ0$, is replaced by the hypothesis $\E\|\hat T\|^2<\infty$. Step~B
uses three properties of the distortion, and each holds for $\ell_{\mathcal C}$.
(1) The finite-alphabet code of \cite[Lemma~A]{part1} needs $d_k$ bounded on the finite
alphabets. It is, because $\ell_{\mathcal C}(t,\hat x)\le M_1^2\|t-\hat x\|^2$, every cell
has positive probability, and $T$ has finite second moments.
(2) The code's consumer loss equals its surrogate distortion: its reproductions are
functions of $T_k^n$, and given $T_k^n$ the letter $T_i$ depends on $T_{k,i}$ alone.
(3) The surrogate distortion of the channel tends to $\E\,\ell_{\mathcal C}(T,\hat T)$ as
$k\to\infty$. Write
$d_k(t,\hat x)=\E[\|g(T)\|^2\mid T_k=t]-2\langle\E[g(T)\mid T_k=t],g(\hat x)\rangle+\|g(\hat x)\|^2$.
The surrogate law of the pair $(T_k,\hat T_k)$, with $\hat T_k=q_k(\hat T)$, is its true
law, so the surrogate distortion is
\[
\E\|g(T)\|^2-2\,\E\bigl\langle\E[g(T)\mid T_k],\,g(\hat T_k)\bigr\rangle+\E\|g(\hat T_k)\|^2 .
\]
Here $g(T)$ is square-integrable, because $g$ is Lipschitz and $T$ is Gaussian;
$\E[g(T)\mid T_k]\to g(T)$ in $L^2$ by martingale convergence, because the partitions
of the quantizers are nested and generate the Borel $\sigma$-field; and
$g(\hat T_k)\to g(\hat T)$ in $L^2$, because $\hat T_k\to\hat T$ in $L^2$
\cite[proof of Theorem~3]{part1} and $g$ is Lipschitz. So the inner product and the last
term converge, and the surrogate distortion tends to
$\E\|g(T)-g(\hat T)\|^2=\E\,\ell_{\mathcal C}(T,\hat T)$. The rest of the proof is unchanged.
For the test channel, $I(T;\hat T)=r(\Delta)$ and $I(T;\hat T\mid S)=\ell(\Delta)$
\cite[Proposition~2]{part1}.
\end{proof}

The code is built for the consumer's surrogate distortion, so it depends on
$\mathcal C$, not only on $Q$; Remark~\ref{rem:open} returns to this. Its errors need not
be Gaussian: what transfers is the consumer loss, through the convergence in~(3).

\begin{corollary}[Learn and certify without the lower sandwich]\label{cor:cert}
In the setting of Theorem~\ref{thm:second}, suppose $Q\succ0$ and the encoder knows
$\Sigma$, the constants $M_1,M_2,M_3$, and a $\nup$ with $\PC\preceq(1+\nup)Q$ (for instance
$\nup=M_1^2/\lambda-1$ for an a priori $\lambda\le\lmin(Q)$), and let $\varepsilon_\pm$
be those of Lemma~\ref{lem:est} with $\num=1$, $\emm<1$. Let $\bar\rho(D)$ be $\rho_b$
with $s$ replaced by $D/\lmin(\hat Q)$, and suppose $\bar\rho(D)<D$. On the event of
Lemma~\ref{lem:est}, with $\hat D=(1-\emm)(D-\bar\rho(D))$:
\begin{itemize}
\item[(a)] \emph{Test channels.} Every Gaussian test channel with
$\tr(\hat Q\hat\Delta)\le\hat D$ has consumer loss at most $D$; in particular, for every
$\alpha$, so does the one that attains $V_{\hat Q}(\alpha,\hat D)$, and
\[
V_{\hat Q}(\alpha,\hat D)\le V_Q(\alpha,D)+\bar\tau,\qquad
\bar\tau=\frac p2\log\frac{1+\epp}{(1-\emm)(1-\bar\rho(D)/D)}=\tau_\varepsilon+O(D).
\]
\item[(b)] \emph{Codes.} Every pair of $\cR_{\hat Q}(\hat D)$ is achievable at consumer
loss $D$: the operational region of the consumer loss at $D$ contains
$\cR_{\hat Q}(\hat D)$.
\end{itemize}
\end{corollary}

\begin{proof}
(a) Such a $\hat\Delta$ has $\tr(\hat Q\hat\Delta)\le\hat D\le D$, so
$\tr\hat\Delta\le D/\lmin(\hat Q)$ and $\rho(\hat\Delta)\le\rho_b(\hat\Delta)\le\bar\rho(D)$,
since $\rho_b$ is increasing in $s$. On the event,
$\tr(Q\hat\Delta)\le\tr(\hat Q\hat\Delta)/(1-\emm)$, and Theorem~\ref{thm:second} gives
$\E\ell_{\mathcal C}\le\hat D/(1-\emm)+\bar\rho(D)=D$. For the cost,
$V_{\hat Q}(\alpha,\hat D)\le V_{(1+\epp)Q}(\alpha,\hat D)
=V_Q(\alpha,\hat D/(1+\epp))\le V_Q(\alpha,D)+\bar\tau$ by (P2), (P1), and (P3) with
$c=D(1+\epp)/\hat D\ge1$.
(b) Every pair of $\cR_{\hat Q}(\hat D)$ lies above the rate and leakage of a test
channel with $\tr(\hat Q\hat\Delta)\le\hat D$ \cite[Proposition~2, Theorem~3]{part1},
whose consumer loss is at most $D$ by (a). By Theorem~\ref{thm:codes}, for every
$\epsilon>0$ there are codes within $\epsilon$ of that pair and of the loss $D$; a
diagonal sequence as $\epsilon\downarrow0$ achieves the pair at consumer loss $D$, and
achievable pairs are closed upward.
\end{proof}

$V_Q(\alpha,D)$ here is the series' frontier for the true weight, which the oracle's
test channels attain at consumer loss $D+O(D^2)$; without (A1a) it is a benchmark and not a
lower bound. The certificate is conservative: in the checks (C6) the certified designs
have consumer loss between $0.53D$ and $0.90D$, mostly because without a lower constant
$\num=1$ makes $\emm$ large. The encoder needs $\Sigma$ for the design but not for
$\bar\rho$.

\begin{remark}[What remains open for codes]\label{rem:open}
Three things. First, one code for a class: the codes of Theorem~\ref{thm:codes} are
built for the consumer's surrogate distortion, so (A3) covers the test-channel design of
Corollary~\ref{cor:cert}(a), with the derivative bounds among the shared constants, but
not the codes of (b), which need more of $\mathcal C$ than $\hat Q$. A single code, chosen from $\hat Q$ and the constants, that serves every
consumer of the class would follow from a form of \cite[Lemma~A]{part1} that holds for
all bounded distortions at once; a joint-typicality construction suggests such a form,
and this note has not checked it. Second, the process settings
\cite{part2,part3,part4}, whose coding theorems pass to the limit differently. Third, a
converse: without the sandwich of the next section, $V_Q$ is a benchmark only. The
route through \cite[App.~A]{part1} is the author's suggestion.
\end{remark}

\section{Consumers under a quadratic sandwich}\label{sec:sandwich}

For consumers whose loss is pinned between two multiples of one quadratic form, a
converse holds in every setting of the series, and the series' own codes for the
learned weight serve the consumer, with no need to know more of it.

\textbf{Assumption (A1), relative quadratic sandwich.} $Q\succ0$, and there
are $\num,\nup\in[0,1)$, known to the encoder in advance, such that
\begin{itemize}
\item[(A1a)] for every $x$ in the support of the source's symbol laws and every $\hat x$,
\begin{equation}
(1-\num)(\hat x-x)^{\mathsf T}Q(\hat x-x)\ \le\ \ell_{\mathcal C}(x,\hat x)\ \le\ (1+\nup)(\hat x-x)^{\mathsf T}Q(\hat x-x);
\label{eq:A1}
\end{equation}
\item[(A1b)] $(1-\num)Q\preceq\PC(x)\preceq(1+\nup)Q$ for every $x$ in the support of
$\pi$ at which $g$ is differentiable.
\end{itemize}
The bounds must come from prior knowledge of the consumer, such as Lipschitz
and monotonicity constants of the layers of $g$: (A1a) is a statement about all pairs
of points, and no finite set of probes certifies it. For a linear consumer
$\mathcal C(x)=Ax$ with $G^{1/2}A$ of full column rank, $\num=\nup=0$ and
$Q=A^{\mathsf T}GA$. (A1b) follows from (A1a) at points of differentiability by
letting $\hat x\to x$, when $\pi$ is supported in the source's support. The upper half
of (A1a) makes $g$ Lipschitz, with constant $((1+\nup)\lmax(Q))^{1/2}$, on the
interior of the source's support; there $g$ is differentiable almost everywhere
(Rademacher), so for $\pi$ absolutely continuous and supported there the read metric
exists at almost every probe point. Nothing here makes $\mathcal C$ itself Lipschitz or
differentiable when $G$ is singular (Remark~\ref{rem:restrictive}). Only
Lemma~\ref{lem:est} uses (A1b), and it needs no lower constant.

\begin{proposition}[The price of a sandwich]\label{prop:sand}
Let $W\succ0$ and $0<a\le b$.
\begin{itemize}
\item[(i)] If $\ell_{\mathcal C}(x,\hat x)\ge a(\hat x-x)^{\mathsf T}W(\hat x-x)$ for every $x$
in the support of the source's symbol laws and every $\hat x$, then every sequence of
codes, with or without private randomness, whose consumer loss per symbol has limit
superior at most $D$ has every limit point of its pairs of rate and leakage in
$\cR_W(D/a)$.
\item[(ii)] If $\ell_{\mathcal C}(x,\hat x)\le b(\hat x-x)^{\mathsf T}W(\hat x-x)$ on the same
set, then the codes given by the setting's achievability theorem for the weight $W$
at the distortion $D/b$ have consumer loss per symbol with limit superior at most $D$,
and they approach every pair of
\[
\cR_W(D/b)\ \supseteq\ \cR_W(D/a)+\tfrac p2\log\tfrac ba\,(1,1).
\]
\end{itemize}
So a loss known only to lie between $a$ and $b$ times a quadratic form costs at most
$\tfrac p2\log(b/a)$ bits per symbol over the outer bound, in rate and in leakage at
once.
\end{proposition}

\begin{proof}
(i) Applied to each symbol, the lower inequality gives $W$-distortion at most $D'/a$
for every $D'>D$ eventually. The setting's converse holds for codes with or without
private randomness \cite[Sec.~II]{part4} (for a Gaussian vector by
\cite[Proposition~2]{part1}, which applies to any variable from which the
reproductions are computed), so every limit point lies in $\cR_W(D'/a)$ for every
$D'>D$, hence in $\cR_W(D/a)$ by the continuity noted in the setting.
(ii) Applied to each symbol, the upper inequality bounds the consumer loss by $b$
times the $W$-distortion, whose limit superior is at most $D/b$ for the setting's
codes. The inclusion is (P3) with $c=b/a$.
\end{proof}

\begin{theorem}[Learn then serve]\label{thm:T1}
Let $D>0$ and suppose (A1), (A2), (A3), and $\emm<1$. Let $\mathcal E$ be the event of
Lemma~\ref{lem:est}, of probability at least $1-\delta$.
\begin{itemize}
\item[(a)] \emph{Converse, for any learner.} Every sequence of codes, with or without
private randomness, whose consumer loss per symbol has limit superior at most $D$ has
every limit point of its pairs in $\cR_Q(D/(1-\num))$.
\item[(b)] \emph{Achievability, learned.} On $\mathcal E$, the codes given by the
achievability theorem of the setting for the weight $\hat Q$ at the distortion
$\hat D=D(1-\emm)/(1+\nup)$ have consumer loss per symbol with limit superior at most
$D$, and they achieve every pair of
\[
\cR_Q\bigl(D/(1-\num)\bigr)+\tau\,(1,1),\qquad
\tau=\tau_\nu+\tau_\varepsilon,\quad
\tau_\nu=\frac p2\log\frac{1+\nup}{1-\num},\quad
\tau_\varepsilon=\frac p2\log\frac{1+\epp}{1-\emm}.
\]
The encoder uses $\hat Q$, $\nup$, $\emm$, and $\epp$ (so $\nu_\pm$, $m$, and $\delta$),
and the decoder holds the same $\hat Q$ by (A3); the source model and the holder's model
are those the setting's achievability theorem assumes known.
\item[(c)] \emph{The gap.} On $\mathcal E$ the learned code is within $\tau$ bits per
symbol, in rate and in leakage at once, of the outer bound of (a), hence of the best
any code can do at consumer loss $D$. The learning part
$\tau_\varepsilon=(\epp+\emm)p/(2\ln2)+O(\varepsilon^2)$ is sharp, for Gaussian vector
sources, for an encoder that is given $\hat Q$ and $\varepsilon_\pm$ but not the probes
(Proposition~\ref{prop:sharp}); whether the probe data themselves allow less is a
question of statistical lower bounds that this note does not treat. The approximation
part $\tau_\nu$ remains with exact knowledge of $Q$ and bounds the price of the
quadratic approximation from above. The test channels of the series do not pay its
factor $1+\nup$ as $D\to0$ (Theorem~\ref{thm:second}), and in the memoryless setting,
for a consumer whose weighted map has three bounded derivatives, neither do codes built
for that consumer (Corollary~\ref{cor:cert}, Remark~\ref{rem:open}); its factor
$1-\num$ is the converse's.
\end{itemize}
\end{theorem}

\begin{proof}
(a) is Proposition~\ref{prop:sand}(i) with $W=Q$ and $a=1-\num$.
(b) On $\mathcal E$, $Q\preceq\hat Q/(1-\emm)$ and $Q\succeq\hat Q/(1+\epp)$, so (A1a)
gives the sandwich of Proposition~\ref{prop:sand} for the weight $\hat Q$ with
\[
a=\frac{1-\num}{1+\epp},\qquad b=\frac{1+\nup}{1-\emm},\qquad \frac ba=\frac{(1+\nup)(1+\epp)}{(1-\num)(1-\emm)},
\]
and $D/b=\hat D$. Proposition~\ref{prop:sand}(ii) bounds the loss and gives
$\cR_{\hat Q}(\hat D)\supseteq\cR_{\hat Q}(D/a)+\tau(1,1)$, since
$\tfrac p2\log(b/a)=\tau$. Finally $\hat Q\preceq(1+\epp)Q$, so by (P2) and (P1)
$\cR_{\hat Q}(D/a)\supseteq\cR_{(1+\epp)Q}(D/a)=\cR_Q(D/(a(1+\epp)))=\cR_Q(D/(1-\num))$.
(c) is (a) and (b) together.
\end{proof}

The two sandwiches compose: the approximation constants and the estimation
constants multiply into $a$ and $b$, and $\tau$ is the logarithm of their ratio.

\begin{proposition}[The learning price is sharp]\label{prop:sharp}
In the setting of \cite{part1}, with $T\sim N(0,\Sigma)$, the decoder reproducing $T$,
and a holder $S=HT+U$ with information matrix $\mathsf J=H^{\mathsf T}\Sigma_U^{-1}H$,
let $W\succ0$ and $0<a\le b$. For every $\alpha\in[0,1]$ and
$0<D\le a\,\lmin(W)\lmin(\Sigma)$,
\[
\frac p2\log\frac ba-(1-\alpha)\frac{\lmax(\mathsf J)\,D\,(1/a-1/b)}{2\ln2\;\lmin(W)}
\ \le\ V_W(\alpha,D/b)-V_W(\alpha,D/a)\ \le\ \frac p2\log\frac ba .
\]
Consequently, consider an encoder that is given $\hat Q$ and $\varepsilon_\pm$ but not
the probe data, and must keep the consumer loss at most $D$ for every consumer whose
weight $Q$ is consistent with $(1-\emm)Q\preceq\hat Q\preceq(1+\epp)Q$. Among these are
the linear consumers with weights $Q_1=\hat Q/(1-\emm)$ and $Q_2=\hat Q/(1+\epp)$. Safety
for the first forces a weighted cost of at least $V_{\hat Q}(\alpha,D(1-\emm))$, while
the best code for the second costs $V_{\hat Q}(\alpha,D(1+\epp))$. So, for
$D\le\lmin(\hat Q)\lmin(\Sigma)/(1+\epp)$, the regret against the second is at least
\[
\tau_\varepsilon-(1-\alpha)\frac{\lmax(\mathsf J)\,D\,(\epp+\emm)}{2\ln2\;\lmin(\hat Q)},
\]
which tends to $\tau_\varepsilon$ as $D\to0$ and equals it in rate ($\alpha=1$). The
design for $\hat Q$ at $D(1-\emm)$ is safe for every linear consumer consistent with the
event and has regret at most $\tau_\varepsilon$ by (P3), so it is minimax in this limit.
Two limits: for linear consumers a learner that sees the probes recovers $Q$ exactly,
so this is a statement about the confidence statement, not a statistical lower bound;
and for $\nu_\pm>0$ it shows that the total gap is at least $\tau_\varepsilon-O(D)$, not
that $\tau_\varepsilon$ is paid on top of $\tau_\nu$.
\end{proposition}

\begin{proof}
The upper bound is (P3). For the lower bound let $\Delta'$ attain $V_W(\alpha,D/b)$
and put $c=b/a$. From $\tr(W\Delta')\le D/b$, $\tr\Delta'\le D/(b\lmin(W))$, so
$\|c\Delta'\|\le D/(a\lmin(W))\le\lmin(\Sigma)$ and $c\Delta'$ is feasible at the
distortion $D/a$. Its rate is $r(\Delta')-\tfrac p2\log c$. Its leakage is
\[
\ell(c\Delta')=\tfrac12\log\det(c^{-1}\Delta'^{-1}+\mathsf J)-\tfrac12\log\det(\Sigma^{-1}+\mathsf J)
=\ell(\Delta')-\tfrac p2\log c+\tfrac12\log\frac{\det(I+c\Delta'\mathsf J)}{\det(I+\Delta'\mathsf J)}.
\]
The eigenvalues $\mu_i$ of $\Delta'\mathsf J$ are those of
$\Delta'^{1/2}\mathsf J\Delta'^{1/2}\succeq0$, and
$\ln\frac{1+c\mu}{1+\mu}\le(c-1)\mu$, so the last term is at most
$(c-1)\tr(\Delta'\mathsf J)/(2\ln2)\le(c-1)\lmax(\mathsf J)D/(2\ln2\,b\lmin(W))$, and
$(c-1)/b=1/a-1/b$. So $V_W(\alpha,D/a)\le\alpha r(c\Delta')+(1-\alpha)\ell(c\Delta')$ is
at most $V_W(\alpha,D/b)-\tfrac p2\log c+(1-\alpha)\lmax(\mathsf J)D(1/a-1/b)/(2\ln2\,\lmin(W))$.
For the consequence: the first consumer's loss is $\delta^{\mathsf T}\hat Q\delta/(1-\emm)$,
so Proposition~\ref{prop:sand}(i) with $W=\hat Q$ and $a=1/(1-\emm)$ bounds the cost
below by $V_{\hat Q}(\alpha,D(1-\emm))$; the second's region is
$\cR_{Q_2}(D)=\cR_{\hat Q}(D(1+\epp))$ by (P1); apply the bound with $W=\hat Q$,
$a=1/(1+\epp)$, $b=1/(1-\emm)$, for which $1/a-1/b=\epp+\emm$.
\end{proof}

\section{Scope}

\begin{remark}[How restrictive (A1) is]\label{rem:restrictive}
For a Gaussian source the support is all of $\R^p$, so (A1a) is global: its lower half
makes the weighted map $g=G^{1/2}\mathcal C$ injective with a Lipschitz inverse on its
image, and its upper half makes $g$ Lipschitz, so $g$ is bi-Lipschitz on $\R^p$. When
$G\succ0$ so is $\mathcal C$. When $G$ is only semidefinite, $\mathcal C$ need not be
Lipschitz: $\mathcal C(x)=(x,x^2)$ on $\R$ with $G=\mathrm{diag}(1,0)$ has
$\ell_{\mathcal C}(x,\hat x)=(\hat x-x)^2$, so (A1a) holds with $Q=1$ and $\nu_\pm=0$,
while $x^2$ is not globally Lipschitz. The restrictions are on $g$. A bi-Lipschitz $g$
excludes a square ReLU layer, which is not injective, an output metric of rank below $p$
(in particular a map to fewer than $p$ outputs), a softmax, polynomials of degree two or
more (a globally Lipschitz polynomial is affine), and quantizers; for $G\succ0$ these
exclusions are of $\mathcal C$ itself, and the part of $\mathcal C$ in the kernel of $G$ is
free. An expanding ReLU layer can satisfy (A1): applied componentwise,
$g(x)=(\mathrm{ReLU}(x),\mathrm{ReLU}(-x))$ has $\PC=I$, and (A1a) holds with $Q=I$,
$\num=1/2$, and $\nup=0$, so $\tau_\nu=p/2$ bits per symbol. It admits leaky-ReLU layers with invertible weights, invertible
residual blocks $x+f(x)$ with $\mathrm{Lip}(f)<1$, and monotone maps with derivative in
$[a,b]$, $a>0$. Even then $\tau_\nu$ is small only for nearly linear consumers: a
componentwise leaky ReLU of slope $a$ under a symmetric $\pi$ has
$\num=\nup=(1-a^2)/(1+a^2)$, so $\tau_\nu=p\log(1/a)$ bits per symbol, $6.6p$ at
$a=0.01$ and $2.3p$ at $a=0.2$. (A1a) is what the converse and the operational
achievability of Theorem~\ref{thm:T1} need in every setting of the series. For the test
channels, Theorem~\ref{thm:second} needs only three bounded derivatives of $g$,
Theorem~\ref{thm:codes} only that $g$ be Lipschitz, and Corollary~\ref{cor:cert} also
$Q\succ0$; a leaky ReLU is not $C^2$, and a smoothed one has $M_2$ and $M_3$ of the
order of the inverse smoothing width and its square.
\end{remark}

\begin{remark}[Read operators of lower rank]
(A1) as stated needs $Q\succ0$, but its form allows a singular $Q$ with a
kernel that does not move: for $\mathcal C(x)=f(Wx)$, $W$ of rank $r<p$, and $f'$
in $[a,b]$ with $a>0$, both the loss and $\delta^{\mathsf T}Q\delta$ vanish on
$\ker W$ and the sandwich holds. (A1b) then forces
$\ker\PC(x)=\ker Q$ at every probe point; Lemma~\ref{lem:est} works on the
range of $Q$ with $Q^{+1/2}$ and $L=\ln(2r/\delta)$; (P1)--(P3) hold with
$r$ in place of $p$ for the frontier of the reproduction of $Fx$,
$F^{\mathsf T}F$ spanning the range; and \cite{part1,part2,part3} allow the
decoder to reproduce such a linear function of the source, while
\cite{part4} does not. So Theorem~\ref{thm:T1} extends to consumers of lower rank
with a fixed kernel in the settings of \cite{part1,part2,part3}. The open
cases are a kernel that moves with $x$ and a range that is itself estimated
with error. For the rate alone, the foundational manuscript's commission tax
\cite{otieee} prices mis-weighting on an exactly known range (its hypothesis
$m\hat P\preceq P\preceq M\hat P$ forces equal ranges), with the sharp constant
$\tfrac r2\log(M/m)$, and is the rate-only analog of Proposition~\ref{prop:sand};
its omission floor is a floor on the distortion, not a rate price, for a
code blind to a read direction. Principal angles enter there only through
the overlap discussion that follows it.
\end{remark}

\begin{remark}[Losses that are quadratic only near the diagonal]
A consumer whose loss is quadratic near $\hat x=x$ but not globally fails (A1a) even
when its read metric is constant: $\mathcal C(x)=(\cos x,\sin x)$ on $\R$ has
$\PC\equiv1$, yet its loss saturates. For arbitrary codes a bound then needs a tail
term $\E[\ell_{\mathcal C}\mathbf 1\{\|\delta\|>r_0\}]$ that depends on the code's error
law beyond its second moment, which the series does not control; for the Gaussian
test channels Theorem~\ref{thm:second} supplies it, and Theorem~\ref{thm:codes} carries
it to codes in the memoryless setting. A softmax fails (A1a) for a different reason: its
read metric tends to zero at saturated points, so it is not bounded below by a multiple
of $Q$; and the softmax's Jacobian annihilates the all-ones direction of its logits, so
a read-out $\mathrm{softmax}(Wx)$ has a kernel exactly when the range of $W$ contains
that direction. Theorem~\ref{thm:second} applies to it as to any consumer with three
bounded derivatives. Corollary~\ref{cor:cert} also needs $Q\succ0$, which holds for the
registration draft's read-out of two inputs into three outputs (the all-ones vector is
outside the range of its $W$); the audit (VI-23) found $\lmin(Q)\approx2\times10^{-4}$
there, so a certificate would be weak. The leading-order statement of
\cite[Prop.~1]{otieee}, with remainder $o(\varepsilon^2)$ and the condition that the
error's second moment be uncorrelated with $\PC(x)$, is the local counterpart; for
the test channel that condition holds up to the term $\E u^{\mathsf T}(\PC(T)-Q)u$
of the proof.
\end{remark}

\begin{remark}[Consumers that read across symbols]
The series' distortion here is a sum over symbols with one weight. A
consumer that reads several symbols jointly, such as attention over the
tokens of a context, has a read metric over a window, and its trace
statistic can misorder codes whose errors correlate across symbols (the
ledger's NEG-9). For sources made of independent blocks, a block can be
taken as one symbol of a larger vector, and \cite{part1} covers any weight
on it; for processes whose blocks are dependent, weights with cross-symbol
structure are outside \cite{part2,part3,part4}.
\end{remark}

\begin{remark}[Noisy consumers]\label{rem:fd}
Proposition~\ref{prop:fd} assumes noiseless queries. Output noise of variance
$\sigma^2$ makes the estimated Jacobian $J+\mathsf N$ with entries of $\mathsf N$ of
variance of order $\sigma^2/h^2$, and then
$\E[(J+\mathsf N)^{\mathsf T}G(J+\mathsf N)]=J^{\mathsf T}GJ+\E[\mathsf N^{\mathsf T}G\mathsf N]$:
a bias of order $\sigma^2/h^2$ that does not shrink with $m$, against the bias of order
$h^2$ of Proposition~\ref{prop:fd}. It is removed by taking two independent difference
Jacobians $\hat J_1,\hat J_2$ at each point and using
$\tfrac12(\hat J_1^{\mathsf T}G\hat J_2+\hat J_2^{\mathsf T}G\hat J_1)$, at twice the
queries, or by subtracting the bias when $\sigma$ is known. With Gaussian noise the
summands are not bounded, and Lemma~\ref{lem:est} needs a sub-exponential form of
matrix Bernstein; the constants are not worked out here. Rounding in floating point is
such a noise, so $h$ cannot be taken to zero in practice.
\end{remark}

\begin{remark}[The probe distribution]
$Q$ is the average of $\PC$ under $\pi$. Under (A1) the weight of any probe
distribution supported in the source's support lies between $1-\num$ and
$1+\nup$ times $Q$, so two such weights are within the factors
$(1-\num)/(1+\nup)$ and $(1+\nup)/(1-\num)$ of each other; the choice of $\pi$
changes the constants of Theorem~\ref{thm:T1} and not its form.
Theorem~\ref{thm:second} singles out the source law: only for it is the leading
term $\tr(Q\Delta)$. For the sources of \cite{part3,part4}, whose symbol laws change
with time, (A1) is a condition on every symbol law at once.
\end{remark}

\begin{remark}[What the probability means]
The event of Lemma~\ref{lem:est} is a confidence statement about the learned
operator. It is not the false-clear rate of a monitor that certifies a
running description; that needs its own statistical argument and
measurement.
\end{remark}

\begin{remark}[The leakage as work]
The theorems bound the conditional information that is the leakage. Its
reading as the minimal work to reset a record, under the holder's
knowledge, uses the physical memory and reset model of the Landauer
manuscript \cite{landauer}, which is separate from the consumer's model.
\end{remark}

\begin{thebibliography}{9}
\bibitem{part1} A. H. Bond, ``Rate, leakage, and distortion for Gaussian
sources: Water-filling on a tilted weight,'' manuscript IT-26-1911, under review at
\emph{IEEE Trans. Inf. Theory}, 2026 (\texttt{paper/tit-rate-leakage}). Theorem~3 is
its coding theorem, proved in its Appendix~A with Lemma~A there.
\bibitem{part2} A. H. Bond, ``Rate, leakage, and distortion for stationary
Gaussian sources: Water-filling on a tilted spectrum,'' manuscript, 2026,
tag \texttt{tit-stationary-1.4}.
\bibitem{part3} A. H. Bond, ``Rate, leakage, and distortion beyond
stationary Gaussian sources: Water-filling over time and frequency,''
manuscript, 2026, tag \texttt{tit-beyond-stationary-1.3}.
\bibitem{part4} A. H. Bond, ``Rate, leakage, and distortion for unstable and
rapidly varying Gaussian sources,'' manuscript, 2026
(\texttt{paper/tit-unstable-ergodic}, merged as pull request 26).
\bibitem{otieee} A. H. Bond, ``Observation theory'' (the foundational
manuscript), \texttt{paper/ieee/observation-theory-ieee.tex}: Proposition
``Consumer-projected distortion'', Proposition ``Recovery of the read
subspace'', and the theorems ``Commission tax'' and ``Omission floor''.
\bibitem{landauer} A. H. Bond, ``Consumer-relative Landauer'' (Paper V),
\texttt{paper/consumer-relative-landauer.tex}.
\bibitem{tropp2012} J. A. Tropp, ``User-friendly tail bounds for sums of
random matrices,'' \emph{Found. Comput. Math.}, vol.~12, no.~4,
pp.~389--434, 2012.
% Theorem 1.4 (matrix Bernstein, one-sided hypothesis lambda_max(X_k) <= R) checked
% against arXiv:1004.4389v7 on 2026-11-17; the journal numbering was not checked.
\bibitem{bhatiadavis} R. Bhatia and C. Davis, ``A better bound on the
variance,'' \emph{Amer. Math. Monthly}, vol.~107, no.~4, pp.~353--357, 2000.
% Used for attribution of the scalar inequality only; the matrix form is proved in the text.
\bibitem{linderzamir} T. Linder and R. Zamir, ``High-resolution source coding
for non-difference distortion measures: The rate-distortion function,''
\emph{IEEE Trans. Inf. Theory}, vol.~45, no.~2, pp.~533--547, 1999.
% Bibliographic data from DBLP; the paper was not read at first hand.
\end{thebibliography}

\end{document}
